\section{統計環境Rを使ってみる}
\label{s04_usingR}

\subsection{授業内容}
\subsubsection{科目の中でのこのコマの位置づけ}
\begin{tabular}[t]{ccccccc}
    記述統計[2] & 確率[0] & モデル[0] & 推測・推定[0] & 意思決定[0] & 実践[2]
\end{tabular}
\subsubsection{概要}
統計処理を行うに当たっては，パソコンの利用が不可欠である。また，ExcelやNumbersなどの表計算でも記述統計量やグラフの作図はできるが，それ以上の「複数の変数を扱った多変量解析」や「群間の平均値の差を細かく検証する」といった心理統計的な作業には不十分でる。そこで，より専門的なツールとして統計環境Rを使うことになる。統計環境Rと，総合開発環境であるRStudioの基本的な使い方，結果の再現性を担保するための工夫などについても言及する。
\subsubsection{コマ主題細目}
\begin{description}
\item[Rmarkdownによる動的な文書作成]RStudioをRの実行環境として捉えるだけでなく，Rのもつ豊かなパッケージ群とRStudioのもつ統合的な環境とを合わせて活用することで，文書作成を1つの環境の中で統合的に扱うことができるようになる。RmarkdownというRの書式の1つについて理解し，データの処理と報告書の作成を統一された環境で行うことで，複数の環境を跨ぐことによって生じる問題(ex.コピペ汚染)から脱却し，再現性のある実行環境を維持する取り組みについて理解する。
\begin{flushright}
    $\to$ \citet{kosugi2019}のP.171--188。再現可能性も視野に入れたより専門的な参考図書は\citet{高橋201805}。
\end{flushright}

\item[外部データの読み込み]


\item[整然データ]データを分析する際には，データの持っている尺度水準と特定の値に紐づけられた情報が，行と列の指定で一意に特定できることが重要である。このようなデータのことを特に整然データという。整然データを目標としたデータの入力の必要性を理解する。また不適切に入力されたものであっても整然データの形に整形するための技術を理解しなければならない。
\begin{flushright}
    $\to$ 整然データについては\url{https://id.fnshr.info/2017/01/09/tidy-data-intro/}がわかりやすい。
\end{flushright}

\item[ggplot2によるグラフの文法]大量のデータが整然とした形に整えられ，そこから情報を汲み取る時にまず行うのが可視化であり，データの要約統計量による理解である。ここでは先のデータの整然化を経て，理解を進めるグラフを描画するための文法について理解する。Rのggplot2パッケージをもちいて，ヒストグラム，ボックスプロット，散布図の描画，それらの群ごとの色分けや掻き分けなどがスクリプトによって実現できることを理解し，一度描画言語で記述することで再現性を高めることができることを理解する。
\begin{flushright}
    $\to$ \citet{紀ノ定201806}を参照
\end{flushright}

\end{description}
\subsubsection{キーワード}
\begin{itemize}
\item R
\item RStudio
\item 整然(tidy)data
\item ggplot2
\end{itemize}
\subsection{授業情報}
\paragraph{コマの展開方法}
演習；4号館4FのPCルーム1/2に別れて，実際のPCを使う実習を伴った授業を行う
\paragraph{標準シラバスにおける位置づけ}
科目番号5 心理学統計法；2統計に関する基礎的な知識；C/D/E エクセル，R，SPSS入門
\subsubsection{予習・復習課題}
\paragraph{予習}
受講生の中にはまだ，Excelなどの表計算ソフトを実際に利用・活用したことがないものがいるかもしれない。住所録や家計簿といった身の回りの情報であっても，数値化して入力されたものを一覧にし，必要な関連キーワードによってまとめ，特徴を描写するなどの方法は，高度に情報化された現代社会では基本的な素養の1つである。身の回りでのデータの活用例について事前に情報収集し，本時で学ぶ内容を即座に活用できるような側面についていろいろ考えておくことも予習になる。また有償であるソフトウェアは，品質保証と引き換えにその内部構造が企業秘密＝ブラックボックスであり，そのことを問題視するフリーソフトウェアという考え方がある。このことについては\cite{長尾200305}に詳しく，事前知識として一読すると良いだろう。
\paragraph{復習}
RおよびRStudioを使った実習については，大学のPC環境で利用できるのはもちろんであるが，フリーソフトウェアなので個人所有のPCにインストールすることができる。\citet{kosugi2019}などを参考に，手元の環境で実行環境を構築することで，今後の分析やレポート作成に向けた準備にもなる。
授業時間内に課される課題をRmd形式で提出することが求められる。時間内に学んだ技術を用いて，レポートを提出すること。なおこのレポートの提出は前期の単位認定にあたっての必要条件であり，基準に満たないレポートは再提出を求める。前期試験の前にパスすることが，前期末試験を受験するための前提条件であることを承知しておくこと。
